\section{PixArt-$\alpha$: lograr que DiT entienda el lenguaje natural}

Text-to-image sustituye la condición breve de clase por tokens de texto largos. La arquitectura debe resolver a la vez el modelado interno de los patch de la imagen y la alineación texto--imagen. La propuesta de PixArt-$\alpha$ conserva la self-attention de DiT, lee el texto mediante cross-attention y sigue modulando AdaLN con el timestep.

\subsection{Qué módulos nuevos requiere la condición de texto}

Codificar el prompt en un solo vector pierde las relaciones entre palabras; por eso los modelos predominantes conservan la token sequence del texto. Al leer las dos diapositivas siguientes, conviene seguir por separado dónde se codifica el texto, dónde interactúa con la imagen y cómo entra el paso de tiempo en cada capa. Esta sección pasa de la condición de clase de baja dimensión del Original DiT a la condición en lenguaje natural; el problema no es solo «conectar un encoder más», sino cómo conservar la estructura a nivel de palabra, controlar la dirección de la interfaz y calcular el costo del pipeline completo.

\lecturefigure{slide-35.jpg}{De class-conditional DiT a text-to-image: codificación del texto, cross-attention y modulación temporal}{35}

Tres preguntas forman la lista de decisiones de diseño: qué text encoder usar; si texto e imagen interactúan mediante cross-attention, concatenación o joint attention; y si el timestep sigue usando AdaLN. Pueden elegirse de forma independiente; por lo tanto, «text DiT» no es una arquitectura única.

\lecturefigure{slide-36.jpg}{Bloque de PixArt-$\alpha$: self-attention, cross-attention, MLP y condición de timestep}{36}

Los noisy image tokens pasan primero por self-attention y después por cross-attention, con los image tokens como query y los text tokens como key/value. La condición del paso de tiempo modula el image stream mediante AdaLN. Así, el texto conserva su estructura de secuencia y el número de tokens de imagen puede diferir del número de tokens de texto.

\lecturefigure{slide-37.jpg}{El codificador de texto Flan-T5-XXL admite prompts largos y aporta una comprensión semántica rica}{37}

Un text encoder grande mejora la capacidad semántica, pero también aumenta la memoria de GPU y la latencia. Si el encoder está congelado, el entrenamiento generativo no actualiza la representación del lenguaje; si se entrena de forma conjunta, el costo es mayor y puede deteriorarse la capacidad lingüística. Al comparar tamaños de modelo, es necesario indicar si se incluyen los parámetros del text encoder.

\begin{knowledgebox}{La dirección de la cross-attention}
Que los tokens de imagen actúen como query y los tokens de texto como key/value significa que cada posición espacial lee activamente el significado de las palabras pertinentes. La longitud de salida sigue siendo igual al número de tokens de imagen; esta capa no actualiza directamente el texto. MMDiT, en cambio, hace que ambas modalidades evolucionen juntas.
\end{knowledgebox}

\subsection{AdaLN compartido y resultados de un modelo compacto}

PixArt conserva la timestep modulation, pero no necesita mantener parámetros independientes para cada block. Un modulador compartido reduce el cómputo y los parámetros, y deja las diferencias entre capas a la representación propia de cada block. Esta sección indaga de dónde proviene esa compacidad: lo que se comparte es el mapeo de la condición y no todos los pesos del Transformer; por eso hay que revisar por separado el número de parámetros, el cómputo de entrenamiento y la calidad final, sin suponer que «compartir» equivale automáticamente a mayor rapidez de extremo a extremo.

\lecturefigure{slide-38.jpg}{AdaLN de PixArt: el timestep genera los parámetros de modulación de cada block}{38}

La diapositiva de código muestra que el time embedding, tras pasar por un MLP, produce scale/shift. Igual que en el Original DiT, AdaLN se encarga de la condición temporal; el texto pasa por cross-attention. Separar ambos tipos de condición evita comprimir un texto largo en un único vector de modulación.

\lecturefigure{slide-39.jpg}{Parámetros de AdaLN compartidos con sesgos ligeros para cada block}{39}

Si todos los block comparten el mismo modulador, los parámetros adicionales por capa pueden bajar de $O(LD^2)$ a algo más cercano a $O(D^2+LD)$, donde $L$ es el número de capas. El ahorro por compartir restringe la capacidad expresiva, pero los experimentos muestran que el costo puede reducirse de forma considerable con una pérdida de calidad pequeña.

\lecturefigure{slide-40.jpg}{PixArt-$\alpha$: resultados de texto a imagen competitivos con un número compacto de parámetros}{40}

La tabla y la gráfica de barras deben leerse considerando a la vez los parámetros, el costo de entrenamiento y las métricas. Un modelo más pequeño no implica un entrenamiento más sencillo: la calidad de las caption de los datos, el curriculum de resolución, el VAE, el optimizer y el dropout pueden contribuir al resultado. El ponente deja explícitamente la training recipe para otro tema.

\teachervoice{00:32:56--00:40:45, el docente divide el aporte de PixArt en dos partes: usar un codificador de texto grande y cross-attention para procesar el lenguaje natural; y seguir usando el AdaLN de timestep, compartiendo parámetros para ahorrar cómputo. Que la arquitectura sea sencilla no significa que el entrenamiento sea «fácil».}

\begin{warningbox}{El criterio para contar parámetros debe ser uniforme}
Al reportar los parámetros del denoiser, es posible que no se incluyan el text encoder congelado ni el VAE; sin embargo, el despliegue de extremo a extremo debe cargarlos. Al comparar «modelos más pequeños», conviene listar a la vez el denoiser, los condition encoders, la memoria de GPU máxima, los FLOPs de entrenamiento y la latencia de muestreo.
\end{warningbox}

\subsection{Resumen de la sección}

PixArt-$\alpha$ extiende DiT al lenguaje natural mediante un encoder de texto y cross-attention, al tiempo que conserva el AdaLN de timestep y comparte los parámetros de modulación. Muestra el potencial de un backbone compacto y también recuerda que es necesario separar la arquitectura, la training recipe y el costo del pipeline completo.
